\documentclass[11pt]{article}
\usepackage[a4paper,margin=0.9in]{geometry}
\usepackage{amsmath,amssymb,amsthm,graphicx,microtype,xcolor,mathtools}
\usepackage[colorlinks=true,linkcolor=blue!50!black,urlcolor=blue!60!black,citecolor=blue!50!black]{hyperref}
\newtheorem{theorem}{Theorem}
\newtheorem{lemma}{Lemma}
\newcommand{\R}{\mathbb R}
\newcommand{\V}{\mathcal V}
\newcommand{\W}{\widehat W}
\newcommand{\nuA}{\nu_\alpha}
\newcommand{\dmA}{m_\alpha}
\title{Weak type $(1,1)$ for variation of the Laguerre heat semigroup\\
\large A full answer to the endpoint question in arXiv:2110.03493}
\author{Autonomous mathematical discovery run \texttt{fa\_banach\_001}, lane 5}
\date{12 August 2026}

\begin{document}
\maketitle

\begin{center}
\fbox{\parbox{0.92\linewidth}{\textbf{Status: candidate full, likely valid.}
The proof is analytic.  A short exact symbolic script only audits the degrees of three explicitly displayed polynomials.}}
\end{center}

\section{Source question and answer}

Betancor and de Le\'on-Contreras study variation operators for Laguerre
polynomial expansions in \cite{source}.  On page 9, while proving their
Poisson endpoint theorem, they state that they do not know the weak type
$(1,1)$ bound for the variation of the transformed Laguerre heat semigroup
$\{\W_t^\alpha\}_{t>0}$; see Figure~\ref{fig:source}.

\begin{figure}[ht]
\centering
\includegraphics[width=0.98\linewidth]{figures/open_problem_crop.png}
\caption{The endpoint question and the heat-kernel formula in the source paper \cite[p.~9]{source}.}
\label{fig:source}
\end{figure}

Let $n\geq1$, $\alpha=(\alpha_1,\ldots,\alpha_n)\in(0,\infty)^n$, and
\[
 d\nuA(x)=2^n\prod_{j=1}^n\frac{x_j^{2\alpha_j+1}e^{-x_j^2}}
 {\Gamma(\alpha_j+1)}\,dx_j .
\]
For $\rho>2$, write
\[
 \V_\rho(\{T_t\}_{t>0})f(x)=
 \sup_{0<t_k<\cdots<t_1}
 \left(\sum_{j=1}^{k-1}|T_{t_j}f(x)-T_{t_{j+1}}f(x)|^\rho\right)^{1/\rho}.
\]

\begin{theorem}\label{thm:main}
For every $\rho>2$ there is $C=C(n,\alpha,\rho)$ such that, for every
$f\in L^1((0,\infty)^n,\nuA)$ and $\lambda>0$,
\[
 \nuA\!\left(\left\{x:\V_\rho(\{\W_t^\alpha\}_{t>0})f(x)>\lambda\right\}\right)
 \leq \frac{C}{\lambda}\|f\|_{L^1(\nuA)}.
\]
Consequently the same estimate holds for the original Laguerre polynomial
heat semigroup on $L^1(\mu_\alpha)$, by the isometry $x\mapsto x^2$ used in
\cite[Section 2.1]{source}.
\end{theorem}

The point is that nearly all endpoint machinery needed for the heat family is
already present inside the source's proof for the subordinated Poisson family.
The missing local calculation is supplied below.

\section{Scalar components and bounded total variation}

Put
\[
 A=n+\sum_{j=1}^n\alpha_j,\qquad
 X=|x|^2,\quad Y=|y|^2,\quad Z=\sum_{j=1}^n x_jy_js_j
\]
and, for $r\in(0,1)$,
\begin{equation}\label{eq:G}
 G_{x,y,s}(r)=(1-r^2)^{-A}
 \exp\!\left[-\frac{Xr^2+Y-2Zr}{1-r^2}\right].
\end{equation}
With $r=e^{-t/2}$, the source's integral representation is
\begin{equation}\label{eq:kernel}
 \W_t^\alpha(x,y)=e^{|y|^2}\int_{(-1,1)^n}
 G_{x,y,s}(e^{-t/2})\Pi_\alpha(s)\,ds,
\end{equation}
where
\[
 \Pi_\alpha(s)=\prod_{j=1}^n
 \frac{\Gamma(\alpha_j+1)}{\Gamma(\alpha_j+1/2)\sqrt\pi}
 (1-s_j^2)^{\alpha_j-1/2}.
\]

We use the elementary fact that if a differentiable real function has at most
$d$ critical points, then its total variation is at most
$2(d+1)$ times its supremum norm.  The same conclusion holds when the
derivative numerator below is identically zero.

\begin{lemma}[Uniform critical-point bounds]\label{lem:critical}
For fixed $x,y,s$, off the kernel diagonal,
\begin{align}
 \operatorname{Var}_{0<r<1}G_{x,y,s}(r)&\leq 8\sup_{0<r<1}G_{x,y,s}(r),\label{eq:tvG}\\
 \operatorname{Var}_{0<r<1}\partial_{x_i}G_{x,y,s}(r)
 &\leq 12\sup_{0<r<1}|\partial_{x_i}G_{x,y,s}(r)|,\label{eq:tvx}\\
 \operatorname{Var}_{0<r<1}\partial_{y_i}G_{x,y,s}(r)
 &\leq 10\sup_{0<r<1}|\partial_{y_i}G_{x,y,s}(r)|.\label{eq:tvy}
\end{align}
These are also the total variations under the monotone parametrization
$r=e^{-t/2}$, $t\in(0,\infty)$.
\end{lemma}

\begin{proof}
Logarithmic differentiation of \eqref{eq:G} gives
\begin{equation}\label{eq:P0}
 \frac{dG}{dr}=\frac{2G}{(1-r^2)^2}P_0(r),\qquad
 P_0(r)=Z+(A-X-Y)r+Zr^2-Ar^3.
\end{equation}
Thus $G'$ has at most three zeros unless $P_0\equiv0$, proving
\eqref{eq:tvG}.  Moreover
\begin{equation}\label{eq:spatial}
 \partial_{x_i}G=-\frac{2(r^2x_i-ry_is_i)}{1-r^2}G,\qquad
 \partial_{y_i}G=-\frac{2(y_i-rx_is_i)}{1-r^2}G.
\end{equation}
After differentiating these expressions in $r$ and multiplying by the positive
factor $(1-r^2)^3/G$, the two derivative numerators are polynomials of degrees
at most five and four, respectively.  Hence they have at most five and four
zeros unless identically zero.  The total-variation observation proves
\eqref{eq:tvx}--\eqref{eq:tvy}.  The exact expanded polynomials and degree
checks are reproduced by \texttt{code/verify\_derivative\_degrees.py}.
\end{proof}

Since $\rho\geq1$, every differentiable scalar function $h$ satisfies
\begin{equation}\label{eq:EbyTV}
 \|h\|_{E_\rho}\leq\int_0^\infty|h'(t)|\,dt.
\end{equation}
Lemma~\ref{lem:critical} turns this crude bound into sharp enough kernel
estimates.

\section{Global part}

We use exactly the source's cutoff.  Let
\[
 q_\pm(x,y,s)=|x|^2+|y|^2\pm2\sum_{j=1}^nx_jy_js_j
\]
and choose a smooth $0\leq\varphi\leq1$ which equals $1$ on
\[
 N_1=\left\{q_-(x,y,s)^{1/2}\leq
 \frac{C_0}{1+|x|+|y|}\right\}
\]
and vanishes outside the corresponding enlarged region $N_2$.  Decompose
$\W_t^\alpha=\W_{t,\mathrm{loc}}^\alpha+
\W_{t,\mathrm{glob}}^\alpha$ by inserting $\varphi$ and $1-\varphi$ into
\eqref{eq:kernel}.

The proof of Lemma 2.1 in \cite{source} establishes, before any Poisson-specific
step is used, the following estimate for the heat kernel components:
\begin{equation}\label{eq:globaldom}
 \int_0^\infty|\partial_t\W_{t,\mathrm{glob}}^\alpha f(x)|\,dt
 \leq \mathbb K_\alpha(|f|)(x).
\end{equation}
Indeed, \eqref{eq:tvG} bounds the time integral by a fixed multiple of
$\sup_tG_{x,y,s}(e^{-t/2})$.  On the support of $1-\varphi$, the source proves
\begin{equation}\label{eq:K}
 \sup_tG_{x,y,s}(e^{-t/2})\lesssim K_\alpha(x,y,s):=
 \begin{cases}
 e^{-|y|^2},&Z\leq0,\\[2mm]
 \left(\dfrac{q_+}{q_-}\right)^{A/2}
 \exp\!\left[-\dfrac{|y|^2-|x|^2}{2}
 -\dfrac{\sqrt{q_+q_-}}2\right],&Z>0.
 \end{cases}
\end{equation}
The positive integral operator formed from $K_\alpha(1-\varphi)$ against
$d\dmA(y)=e^{|y|^2}d\nuA(y)$ is $\mathbb K_\alpha$.  By the Sasso estimates
invoked and reproduced in \cite[Lemma 2.1]{source}, it is weak type $(1,1)$
and strong type $(p,p)$ for $1<p<\infty$ with respect to $\nuA$.

Combining \eqref{eq:EbyTV} and \eqref{eq:globaldom},
\begin{equation}\label{eq:globalweak}
 \V_\rho(\{\W_{t,\mathrm{glob}}^\alpha\}_{t>0})f
 \leq\mathbb K_\alpha(|f|),
\end{equation}
so the global variation is weak $(1,1)$ and strong $L^2$.

\section{Local vector-valued Calder\'on--Zygmund estimate}

Let $H_t^\alpha(x,y)$ be the local kernel, now viewed against the doubling
measure $\dmA$:
\[
 H_t^\alpha(x,y)=\int_{(-1,1)^n}
 G_{x,y,s}(e^{-t/2})\varphi(x,y,s)\Pi_\alpha(s)\,ds.
\]
The standard local comparison estimate used in \cite[Lemma 2.2]{source}
implies, on $N_2$,
\begin{align}
 \sup_tG_{x,y,s}(e^{-t/2})&\lesssim q_-(x,y,s)^{-A},\label{eq:sup0}\\
 \sup_t\big(|\nabla_xG_{x,y,s}(e^{-t/2})|+
 |\nabla_yG_{x,y,s}(e^{-t/2})|\big)
 &\lesssim q_-(x,y,s)^{-A-1/2}.\label{eq:sup1}
\end{align}
For completeness, \eqref{eq:sup0} follows by comparing
$q_-(rx,y,s)$ with $q_-(x,y,s)$ and maximizing
$u^{-A}e^{-cq_-/u}$.  For \eqref{eq:sup1}, use
$|rx_i-y_is_i|\leq q_-(rx,y,s)^{1/2}$ (and the analogous $y$ estimate),
absorb the square root into half of the exponential, and maximize once more.

Equations \eqref{eq:EbyTV}, \eqref{eq:tvG}--\eqref{eq:tvy},
\eqref{eq:sup0}--\eqref{eq:sup1}, and
$|\nabla\varphi|\lesssim q_-^{-1/2}$ give
\begin{align}
 \|H^\alpha(x,y)\|_{E_\rho}
 &\lesssim \int_{(-1,1)^n}q_-(x,y,s)^{-A}\Pi_\alpha(s)\,ds
 \lesssim \frac1{\dmA(B(x,|x-y|))},\label{eq:CZ0}\\
 \|\nabla_xH^\alpha(x,y)\|_{E_\rho}+
 \|\nabla_yH^\alpha(x,y)\|_{E_\rho}
 &\lesssim \int_{(-1,1)^n}q_-(x,y,s)^{-A-1/2}\Pi_\alpha(s)\,ds\notag\\
 &\lesssim\frac1{|x-y|\,\dmA(B(x,|x-y|))}.\label{eq:CZ1}
\end{align}
The final integral estimates are precisely the Bessel-kernel estimates used in
\cite[Lemma 2.2]{source}.  Thus the local heat kernel satisfies the same
$E_\rho$-valued Calder\'on--Zygmund bounds as the source's Poisson kernel.

It remains to supply an $L^2$ input.  The family $\{\W_t^\alpha\}$ is a Stein
symmetric diffusion semigroup, so the strong variation theorem of Le Merdy and
Xu \cite[Corollary 6.1]{lemx} yields
\[
 \|\V_\rho(\{\W_t^\alpha\})f\|_{L^2(\nuA)}\lesssim\|f\|_{L^2(\nuA)}.
\]
The global part is strong $L^2$ by \eqref{eq:globalweak}; hence the local part
is strong $L^2$ as well.

Now repeat the source's measure-transfer and vector-valued
Calder\'on--Zygmund argument (its Lemma 2.3 and the paragraph following it).
The local support of $\varphi$ transfers the $L^2$ estimate between $\nuA$ and
the doubling measure $\dmA$.  Bounds \eqref{eq:CZ0}--\eqref{eq:CZ1} then give
weak $(1,1)$ on $\dmA$, uniformly for variation restricted to
$t\in[1/N,N]$.  The same local transfer gives weak $(1,1)$ on $\nuA$.
Letting $N\to\infty$ by monotone convergence proves
\begin{equation}\label{eq:localweak}
 \V_\rho(\{\W_{t,\mathrm{loc}}^\alpha\}_{t>0}):
 L^1(\nuA)\longrightarrow L^{1,\infty}(\nuA).
\end{equation}
Combining \eqref{eq:globalweak} and \eqref{eq:localweak} proves
Theorem~\ref{thm:main}.

\section{Audit, novelty boundary, and limitations}

Run the exact algebra audit from the repository root:
\begin{verbatim}
conda run --no-capture-output -n sandbox python \
  runs/fa_banach_001/solutions/full/\
  2110.03493_laguerre_heat_variation_weak_11/code/\
  verify_derivative_degrees.py
\end{verbatim}
Its captured output is included beside the script.  The computation is not a
numerical experiment: it expands formal derivatives and asserts exact
polynomial identities.

A bounded novelty search on 12 August 2026 covered the run indexes; the exact
arXiv id and question sentence; combinations of ``Laguerre heat semigroup,''
``variation,'' and ``weak type $(1,1)$''; and later endpoint papers citing the
source.  The 2012 Laguerre variation paper \cite{bct} gives strong-$L^p$ and
Hardy-space results in related Laguerre settings.  The 2024 endpoint paper
\cite{endpoint} gives $H^1\to L^1$ and $L^\infty\to\mathrm{BMO}$ estimates.
Neither states this $L^1\to L^{1,\infty}$ heat-variation theorem.  Later
Ornstein--Uhlenbeck work proves analogous Gaussian results but does not cover
arbitrary Laguerre parameters.  No exact prior answer was found; keyword and
citation searches remain incomplete, so novelty requires expert review.

The proof claims only $\rho>2$ and $\alpha\in(0,\infty)^n$, exactly the range
in which the source develops the needed global Sasso estimate.  It makes no
claim at $\rho=2$, for negative Laguerre parameters, or for dimension-free
constants.  The main review-sensitive point is that the local spatial-gradient
total variations, not merely the scalar-kernel variation, are required;
Lemma~\ref{lem:critical} supplies this missing step explicitly.

\paragraph{Human-review recommendation.}
Check the factorization in \eqref{eq:P0} and the two degree bounds following
\eqref{eq:spatial}; then compare \eqref{eq:globaldom}--\eqref{eq:K} and
\eqref{eq:sup0}--\eqref{eq:CZ1} line by line with Lemmas 2.1--2.3 of the source.
Those checks isolate every nonstandard ingredient.

\begin{thebibliography}{9}
\bibitem{source}
J. J. Betancor and M. de Le\'on-Contreras,
\emph{Variation inequalities for Riesz transforms and Poisson semigroups associated with Laguerre polynomial expansions},
arXiv:2110.03493v2 (2022); Anal. Appl. 22 (2024), 1--56.
\url{https://arxiv.org/abs/2110.03493}.

\bibitem{lemx}
C. Le Merdy and Q. Xu,
\emph{Strong $q$-variation inequalities for analytic semigroups},
Ann. Inst. Fourier (Grenoble) 62 (2012), 2069--2097.

\bibitem{sasso04}
E. Sasso,
\emph{Spectral multipliers of Laplace transform type for the Laguerre operator},
Bull. Austral. Math. Soc. 69 (2004), 255--266.

\bibitem{sasso07}
E. Sasso,
\emph{Weak type estimates for Riesz--Laguerre transforms},
Bull. Austral. Math. Soc. 75 (2007), 397--408.

\bibitem{bct}
J. J. Betancor, R. Crescimbeni, and J. L. Torrea,
\emph{Oscillation and variation of the Laguerre heat and Poisson semigroups and Riesz transforms},
Acta Math. Sci. Ser. B 32 (2012), 907--928.

\bibitem{endpoint}
J. J. Betancor, E. Dalmasso, P. Quijano, and R. Scotto,
\emph{Endpoint estimates for harmonic analysis operators associated with Laguerre polynomial expansions},
Math. Nachr. 297 (2024), 2365--2389; arXiv:2210.14394.
\end{thebibliography}

\end{document}
